\chapter{Gauge, Superposition, and the Local Present}
\label{chap:gauge-superposition-present}

Until now the construction has handled one quantity at a time: a number, a duration, a computed
result. This chapter is where it learns to let a measurement vary --- to ask not what a quantity is at
an instant, but how it stands relative to other measurements made elsewhere, and what holds steady when
the comparison is changed. That single shift opens onto three ideas that physics usually presents as
deep and mysterious: a present moment that is only ever local, a gauge freedom that reconciles
observers, and superposition, in which one carrier holds more than one settled answer at once.

A warning is owed at the outset, because this is the first chapter to borrow heavily from the
vocabulary of physics. Every physical word here --- gauge, field, superposition, collapse, the qubit
--- names a finite, combinatorial structure the construction actually builds, with the physics laid
over the top as interpretation. The experiments the chapter leans on are unusually explicit about
this: each one states, in its own text, that no wave dynamics, no metric, and no physical mechanism are
assumed, and that the structure it exhibits is purely informational. We keep that discipline. The
mathematics is earned; the quantum-mechanical reading is an analogy, and a good one, but an analogy.

\section{The local present}
\label{se:the-local-present}

The previous chapter left us with clocks that mint their own order and disagree, so there is already no
global ``now'' to be had. The construction takes the consequence seriously and makes the present a
\emph{local} object. What counts as happening now is a computation carried out in one place, from one
carrier's accumulated memory; the farther away another event is, the farther away in time it
necessarily is, because the only present available is the one assembled locally. A present, in this
construction, is local memory with a deadline: a bounded record of what has just been resolved, good
only here and only for now.

Make the locality concrete with two observers. One stands at the apparatus; the other stands a distance
off and learns of each event only after the news has had time to reach her. Their presents do not
coincide, and the gap is not ignorance to be corrected --- it is structural. The near observer's now
already contains events the far observer's now cannot yet hold, and no amount of careful synchronizing
erases the lag, because the only present either can compute is the one assembled from the records that
have actually reached her. What the construction calls a phenomenon is this richer object: not a bare
event, but an event together with the memory it leaves and the clock-disagreement it provokes between
observers who met it at different local times. A phenomenon is what an event becomes once you admit
there is no single moment at which it happened for everyone.

The deadline is not decoration, and the experiment under Heisenberg's name --- read as a trade-off
rather than a mystery --- shows why. A finite record has a fixed capacity. Spend it refining one group
of outcomes and you have consumed the budget that would have resolved the rest; you cannot sharpen
everything at once, because there is only so much resolving a bounded ledger can do before its deadline
falls. This is the construction's whole account of the uncertainty principle, and it needs no waves:
two measurements are conjugate when they compete for the same finite refinement budget, so that buying
precision in one spends it in the other. The local present is exactly the window inside which that
budget is available; once it closes, the unresolved distinctions are gone, not hidden.

It is worth being concrete about the budget, because it is the whole mechanism. A bounded ledger can
record only so many distinguishable events before its present closes. Two quantities are conjugate when
resolving one to a fine grain requires committing so many of those records that too few remain to
resolve the other. There is no disturbance, no clumsy instrument knocking the system about; there is
only a fixed allowance of distinctions, spent here or spent there but not in both at once. Push all the
budget into pinning a position and the momentum goes unresolved --- not because measuring the position
changed it, but because the records that would have pinned it were never written. The famous inequality
is, on this reading, an accounting identity, the conservation of a finite resolving power, and it is
exact rather than approximate for the same reason a checkbook balances.

\section{The gauge}
\label{se:the-gauge}

If every present is local and every carrier orders events in its own way, then two carriers comparing
notes face a problem: their local orderings need not match, and there is no master frame to appeal to.
The construction's answer is a gauge. A gauge is a transformation that translates one carrier's
local bookkeeping into another's --- and the rule that makes it admissible is the heart of the matter.
A translation between two local frames is permitted only if it leaves the invariant content of the
record unchanged: the number of irreducible refinement steps along any history, the one quantity that
counts genuine events rather than the names given to them, must come out the same on both sides. The
construction tanges exactly the transformations that preserve this count, and funges together all the
local frames that agree on it, treating them as descriptions of one situation.

A concrete relabeling shows what is and is not allowed. Take a history and rename every event in it,
shift the index at which the bookkeeping starts, swap the symbols the carrier uses --- do anything you
like to the description, provided you neither add a step that was not there nor erase one that was. Any
such relabeling is an admissible gauge transformation: the count of genuine refinements is untouched,
every measurable consequence is identical, and the two descriptions are the same situation wearing
different clothes. Now try a transformation that quietly drops a refinement, or invents one. That one is
forbidden, because it changes the invariant count, and a quantity that shifted under it would be an
artifact of the relabeling rather than a fact about the world. This is precisely the gauge freedom of
electromagnetism --- the latitude to change the potential without changing the field --- recovered not
as a postulate but as the bare requirement that physics not depend on how the ledger happens to be
written.

This is the experiment under Maxwell's name, and it reframes gauge freedom as bookkeeping rather than
magic. A transformation that changed the count of irreducible steps would have to introduce or erase
events without recording them, which the construction forbids; so the allowed transformations are
precisely the ones that conserve the event count, and gauge invariance is the demand that physics not
depend on unrecorded relabeling. There is no preferred global frame in this picture and none is needed.
Global structure is assembled from the overlap conditions between many local frames, each complete for
itself --- the field is what you get when you insist that every pair of local observers agree on the
one invariant, and let that agreement knit their separate presents into a whole.

It is worth dwelling on that last phrase, because it inverts the usual picture. One ordinarily imagines
a field as a thing spread through space that local observers sample; here the order of explanation runs
the other way. There is no global field handed down in advance --- there are only many local frames,
each complete for its own carrier, and the constraint that any two of them agree on the invariant count
wherever they overlap. The global object, the field, is nothing but the totality of those overlap
agreements, the structure forced into being by requiring that all the local presents be mutually
consistent. A field, in this construction, is not a substance but a consistency condition with enough
local frames to bind.

\section{Superposition and correlation}
\label{se:superposition-and-correlation}

The last construction of the chapter is the strangest-sounding and, built finitely, the most
demystified. The construction introduces a holder --- think of a jar you can see into but not fully
open --- that carries a state. Usually the state is settled: one definite value, read off cleanly. But
the holder is also allowed to carry a \emph{blend}, two states at once, with neither selected. That
blend is superposition, and the construction is blunt that it is nothing more exotic than a carrier
whose contents are not yet well-ordered --- which, it notes, is the very definition of two things being
simultaneous. A holder in a blend funges two settled states into one bounded object; you can read its
gist without resolving which it is, and for as long as you only read the gist, the blend persists.

What ends a blend is a measurement that cannot coexist with it. Record an event inconsistent with one
of the two branches and the holder must collapse to the other: the construction tanges the surviving
branch, discards the contradicted one, and the blend becomes a single settled state. This is collapse
with no physics in it --- no observer consciousness, no wavefunction, only the rule that the ledger may
not hold a contradiction, enforced the instant an incompatible event is recorded. To carry a blend's
internal orientation --- how the two branches stand relative to each other before collapse --- the
construction needs a quantity that squares to a negative, a marker that is not equal to itself in the
ordinary ordered sense, and that is exactly the role the imaginary unit plays: the bookkeeping symbol
for a rotation between two branches that have not yet been forced to choose.

The imaginary unit earns its place here for a structural reason, not a mystical one. A settled state
sits somewhere on an ordered line; a blend of two states does not, because the two have not been ranked,
and asking which is larger is exactly the question collapse has not yet answered. To track how such an
unranked pair stands --- how far the blend leans toward one branch or the other before any event forces
the issue --- the construction needs a quantity that turns a state by a quarter-turn rather than sliding
it along the line, a marker whose square points backward, opposite to itself. That is the imaginary
unit: not a number on the ordered line at all, but the operator that rotates between the two branches a
blend is holding. It is the same object floating-point arithmetic reaches for when it admits a value not
equal to itself; both are bookkeeping for a state the ordinary ordering cannot place.

\begin{quote}
\emph{A note on the labels.} This section is where the modeling is heaviest, so we mark it plainly. The
holder, the blend, the collapse on contradiction, and the rotation-marker are all finite, combinatorial
constructions. Calling the blend a \emph{quantum superposition}, the collapse a \emph{measurement}, the
two-branch holder a \emph{qubit}, and the rotation-marker \emph{the imaginary unit of a wavefunction}
is interpretation laid over them. The experiments agree: the interferometer and the decohering doublet
both state outright that their content is informational and assume no wave mechanics. What is real here
is the rule that a record may not hold a contradiction; the quantum story is the analogy that rule
turns out to fit.
\end{quote}

The cleanest picture of collapse as a record simply ending is the experiment under Malus's name. Send
light through one polarizing filter and every photon that survives is aligned to its axis; the holder
has been collapsed to a single settled state. Now place a second filter crossed at a right angle to the
first. Nothing gets through. There is no new distinguishable event to record, the count of counts
cannot increase, and the experimental record for that beam simply stops --- the light is no longer
observable, not because it was destroyed but because the apparatus can draw no further distinction from
it. Collapse, termination, and the exhaustion of a record are one phenomenon seen from three angles, and
none of them needs a wave to be doing anything.

Two holders, finally, can be tied together so that resolving one constrains the other --- a correlation
the construction builds by giving the two a shared concept, so that collapsing one forces the matching
branch in its partner. A worked case makes the cluster concrete. Send a single carrier into a device
that splits it down two indistinguishable paths; both paths are admissible, each accumulates its own
consistent history, and the holder genuinely carries two coexisting ledgers at once, exactly as the
Mach--Zehnder experiment describes. Nothing in either arm can say which path is ``real,'' because both
are. Now reunite the paths at a second junction and record a detection. The detection is a strict
update that cannot agree with both histories, so the blend collapses to one --- this is the decohering
doublet of the qubit experiment, a pair of equally admissible refinements pivoted to a single
admissible history the moment an inconsistent event arrives. This correlation is the construction's whole account of entanglement, and it is worth being precise
about what it does and does not permit. Tying two holders to a shared branch means that the instant one
is collapsed the other's matching branch is fixed --- not influenced, fixed, because the two were never
independent; they were one bounded object with two ends. Nothing travels between them when one is
resolved, and nothing can be sent that way: an observer at the far holder sees only its own local
collapse, with no way to tell whether the near holder has been touched yet. What looks from outside like
spooky action at a distance is, from inside, the bare consistency rule again --- a single record may not
hold a contradiction --- applied to a holder whose two ends happen to sit far apart. The correlation is
real and the signal is impossible, and the construction draws both from the same finite discipline that
gave it superposition and collapse.

Superposition, gauge, and the local present
turn out to be three faces of one finite discipline: keep every present local, let observers reconcile
only through what they must agree on, and never let a single record hold two answers once an event has
forced the choice.

\section{What is demonstrated, and what is assumed}
\label{se:demonstrated-assumed-ch5}

Keeping the experimental library's discipline, and with the modeling note above already standing, the
separation here is unusually clean because the experiments draw it themselves.

What is demonstrated is a finite, combinatorial account of three things physics treats as foundational.
There is a local present: a bounded memory with a fixed refinement budget, from which a conjugate
trade-off follows without any wave. There is a gauge: the family of frame-translations that conserve
the count of irreducible events, with global structure assembled from local overlaps and no preferred
frame. And there is a holder that may carry a blend of two states, collapsing to one exactly when an
incompatible event is recorded, with correlated holders built by a shared branch. All of it is finite,
decidable, and free of wave dynamics, metric, or physical mechanism.

\begin{quote}
\emph{A note on the labels.} The chapter's physical names --- \emph{gauge field}, \emph{superposition},
\emph{collapse}, \emph{uncertainty}, \emph{qubit} --- are bridges from these finite structures to the
language of quantum physics, and they are the most ambitious bridges in the book so far. They are
earned to the precise extent that the analogy holds: the construction genuinely exhibits the
informational skeleton of each phenomenon, and genuinely does not exhibit its dynamics. A reader is
entitled to be persuaded that the bookkeeping is the same; a reader is not entitled to conclude that the
construction has derived quantum mechanics, and we do not claim it.
\end{quote}

What is assumed is the interpretive leap itself: that the finite informational structure the
construction builds is the right skeleton for the physical phenomena it is named after. That leap is
the chapter's real content and its real risk, and naming it as an assumption is what lets the next
chapters press further --- assembling, from these pieces, a carrier that can hold a program and run it,
the construction's first complete computer built out of nothing but Facts.
